\section{Audiodatei}\label{audiodatei}

\href{https://sblch-my.sharepoint.com/personal/e037149_sbl_ch/Documents/Transkribierte\%20Datei/WhatsApp\%20Audio\%202026-07-18\%20at\%2018.18.01.mp4}{WhatsApp
Audio 2026-07-18 at 18.18.01.mp4}

\section{Transkript}\label{transkript}

00:00:03 Sprecher 1

nur formell fürs Interview.

00:00:06 Sprecher 1

Das Interview findet auf Hochdeutsch statt.

00:00:10 Sprecher 1

Die IDPA hat Schwerpunkt auf die Fächer Geschichte und Deutsch.

00:00:18 Sprecher 1

Wir werden zwei Biografien verfassen, die Methode bewerten und auch
bewerten wie äußerliche

00:00:30 Sprecher 1

Veränderungen oder dass allgemein, wie die Welt zu dieser Zeit aufgebaut
war, die einzelnen Menschen beeinflusst.

00:00:40 Sprecher 1

Vielen Dank, dass Sie hier sind.

00:00:44 Sprecher 1

Sind Sie einverstanden mit dieser Audioaufnahme?

00:00:48 Sprecher 1

Klar.

00:00:50 Sprecher 2

Sie können auch Du sagen.

00:00:52 Sprecher 1

Okay.

00:00:52 Sprecher 2

Das nehme ich an.

00:00:55 Sprecher 1

Jetzt sind wir per Du.

00:00:59 Sprecher 1

Also, kurze Auflockerung, Sorge.

00:01:04 Sprecher 1

Wie fühlt sich das jetzt an, von Ihrem Neffen interviewt zu werden?

00:01:10 Sprecher 2

Habe ich noch nie gemacht.

00:01:11 Sprecher 2

Von dem her etwas merkwürdig.

00:01:13 Sprecher 2

Aber umgekehrt ist es ja auch ein bekanntes Thema und nicht irgendwie
eine Prüfung sozusagen.

00:01:25 Sprecher 1

Ja, ich kann es nur zurückgeben.

00:01:28 Sprecher 1

Es ist ein interessantes Setting, hatte ich noch nie so.

00:01:31 Sprecher 1

Interviews kenne ich auch nur ziemlich schlichtig.

00:01:37 Sprecher 1

Fangen wir mit dem ersten Punkt an.

00:01:40 Sprecher 1

Wir fangen ziemlich früh an, mit der Kindheit und Ihrem Werdegang in
Basel und mit der ETH.

00:01:48 Sprecher 1

Und eine Grundfrage ist natürlich, wie kam es für Sie zur Informatik?

00:01:55 Sprecher 1

Gab es einen bestimmten Punkt?

00:01:57 Sprecher 1

Hatten Sie eine Person vielleicht?

00:02:00 Sprecher 1

Oder war einfach ein allgemeines Interesse schon da?

00:02:03 Sprecher 2

Ich glaube beides.

00:02:07 Sprecher 2

Damals, ich war ja im Gymnasium, damals hieß das C, Typ C.

00:02:14 Sprecher 2

Das war so Mathewissenschaft.

00:02:17 Sprecher 2

Das hat sich so aufgeteilt.

00:02:19 Sprecher 2

Und wie mein Bruder war ich da im Gymnasium Liestal in der C.

00:02:25 Sprecher 2

das war so ungefähr die Zeit, wo die ersten PCs überhaupt erschienen.

00:02:32 Sprecher 2

Also der Apple 2, vor dem war es, glaube ich, noch der Commodore Pad und
dann das Gimme hatte noch einen alten, so eine Art programmierbaren
Taschenrechner von oder nicht Taschenrechner, der war wirklich so groß
wie eine Schreibmaschine von HP.

00:02:50 Sprecher 2

Also da, da war so ein bisschen was

00:02:53 Sprecher 2

verfügbar und für mich war das einfach interessant, weil es so eine
Kombination ist von Logik und Puzzle und aber auch irgendwas tun, was
dann jemand anderem hilft.

00:03:07 Sprecher 2

Also es ist nicht nur die Mathematik kann ja zum Beispiel sehr trocken
sein, man beweist irgendwas und an der Informatik geht es mehr darum,
etwas zu entwickeln, was dann weiterhilft.

00:03:20 Sprecher 2

oder was eben eine Anwendung ist oder was ein Tool ist oder irgendsowas.

00:03:23 Sprecher 2

Und die Kombination, das sah sehr für mich sehr interessant aus und war
ja auch alles neu.

00:03:31 Sprecher 2

Also nicht nicht irgendwie etwas, wo wo es schon viele Traditionen gab.

00:03:36 Sprecher 1

Man könnte sagen, dein erster Kontaktpunkt mit dem Computer war vor
allem in Gimmi, also im Gymnasium.

00:03:44 Sprecher 1

Und was hast du da genau vielleicht, also

00:03:49 Sprecher 1

Gibt es, war das einfach nur ein Punkt, wo du es zuerst kennengelernt
hast oder hast du sogar schon etwas Spezifisches damit gemacht?

00:03:56 Sprecher 2

Also, ich hab das damals konnte man nicht sehr viel machen.

00:04:00 Sprecher 2

Also zum Beispiel der der Commodore Pet, der der PC am Gimmi war, war
relativ teuer.

00:04:09 Sprecher 2

Das war kostete irgendwie 3000 Franken.

00:04:11 Sprecher 2

Also eigentlich billig damals nicht, aber 3000 Franken war der 8
Kilobyte Speicher und ein

00:04:18 Sprecher 2

Tape Recorder, also ein Kassettengerät, also nicht mal sopidisk als
Speichermedium, das heißt alles war eigentlich super primitiv und super
langsam und man hat dann einfach so ein bisschen eigentlich mehr zum zum
Programmieren lernen irgendwas gemacht, weil du konntest nicht wirklich
ernsthaftes damit tun.

00:04:36 Sprecher 2

Und das andere was ich vergessen habe zu sagen in Sachen Vorbild, es gab
damals schon Niklas Wirth in der Schweiz, der hat ja Pascal erfunden und

00:04:45 Sprecher 2

der war doch schon bekannt, war mal in Zürich an der E.T.H.

00:04:49 Sprecher 2

und also das hat dann bisschen auch noch so eine Schweizer Verbindung
auch gehabt.

00:04:58 Sprecher 1

O.

00:04:59 Sprecher 1

K., ich werde auch mit deiner Schwester oder mit meiner Mutter über die
Kindheit mit dir sprechen.

00:05:05 Sprecher 1

Was denkst du wohl, wird sie also, wenn sie eine Anekdote hätte, was
würde sie mir wohl sagen und wie erinnerst, wie erinnerst du dich daran?

00:05:14 Sprecher 2

Ich weiß genau, was sie sagen wird, und darüber reden wir lieber nicht.

00:05:18 Sprecher 1

Okay.

00:05:20 Sprecher 2

Es involviert Seife.

00:05:23 Sprecher 2

Ja.

00:05:25 Sprecher 2

Okay.

00:05:26 Sprecher 2

Aber sonst glaube ich, also sie war ja schon vier Jahre hintendrein, und
das ist ein großer Unterschied, wenn du selber noch jung bist.

00:05:35 Sprecher 2

Und von dem her, für mich war es eigentlich, also sie sind ja hier
aufgewachsen, war es eigentlich immer ein

00:05:44 Sprecher 2

ein Ort, wo mehrere Kinder rum sind.

00:05:49 Sprecher 2

Also ich, mein Bruder, die Gisela und dann oft noch Schulkollegen oder
sowas.

00:05:56 Sprecher 2

Zum Beispiel die Anita, Freundin von der Gisela, die kam zum Mittagessen
mehrere Jahre.

00:06:02 Sprecher 2

Und Wolfgang, ein jemand, der mit dem Jürgen in die Schule kam, kam auch
zum Mittagessen.

00:06:08 Sprecher 2

Auch mehrere Jahre, weil die so ein bisschen weit weg wohnten und

00:06:12 Sprecher 2

Und damals konnte man an der Schule nichts bekommen.

00:06:17 Sprecher 2

Und von dem her, was ich mich vor allem daran erinnere, ist, dass immer
irgendwo was los war.

00:06:26 Sprecher 1

Meine Mutter hat zu mir gesagt, sie erinnert sich vor allem, dass du
nach der Matura, also jetzt abhand von dem Thema Kindheit, nach der
Matura und dem Militärdienst im IT-Bereich schon Kontaktpunkte hattest,
indem du

00:06:41 Sprecher 1

in einer Versicherung gearbeitet hast.

00:06:43 Sprecher 2

Ja.

00:06:45 Sprecher 1

Was war genau dort dein Kontaktpunkt zu der IT?

00:06:48 Sprecher 2

Das war, also ich habe im Prinzip .

00:06:52 Sprecher 2

Damals hat gerade das Studium .

00:06:55 Sprecher 2

Es gab praktisch keine Möglichkeit, Informatik zu studieren in der
Schweiz damals.

00:07:01 Sprecher 2

Das war also wirklich noch einfach wenig bekannt.

00:07:05 Sprecher 2

Sondern man konnte vielleicht Mathematik studieren und dann sozusagen
Informatik-Schwerpunkt.

00:07:11 Sprecher 2

Und die E.T.H.

00:07:13 Sprecher 2

hatte in dem Jahr wirklich anfangen hätte können, war das wirklich ihr
alles allererstes Jahr, wo es dort Informatik gab.

00:07:22 Sprecher 2

Und ich wollte es aber nicht tun, eigentlich aus 2 Gründen.

00:07:24 Sprecher 2

Erstens mal, ist das allererste Jahr, das ist also so ein One Point O
nicht und ich habe angenommen, es hat Bugs und das wird dann an den an
den Studenten ausgetragen und man hatte auch obligatorischen
Militärdienst

00:07:41 Sprecher 2

Und zusätzlich wollte ich noch was einfach sicher sein, dass die
Informatik auch das Richtige ist für mich.

00:07:50 Sprecher 2

Und dann habe ich eine, es gab damals nicht eine Lehrstelle, aber eine
Programmierstelle bei der Basler angenommen und die haben dich im
Prinzip intern ausgebildet.

00:08:01 Sprecher 2

Und dann hast du auf dem Großrechner programmiert, den PL1 für die
Basler Versicherung.

00:08:09 Sprecher 2

Da war ich dann also

00:08:11 Sprecher 2

nach der guten Schule ein Jahr, ein bisschen mehr als ein Jahr.

00:08:16 Sprecher 1

Wie schon von dir angedeutet, Informatik war als Studienwahl ein
ziemlich junges Fach.

00:08:25 Sprecher 1

Wie wurde es damals eigentlich aufgenommen von deiner Familie?

00:08:30 Sprecher 1

War es ein vernünftiger Weg, den man eingehen sollte oder war das
einfach?

00:08:36 Sprecher 2

Ja, ich glaube, es war eigentlich klar,

00:08:38 Sprecher 2

Also wie gesagt, das war ja im Prinzip alles noch super primitiv, eben
keine PCs, kein Internet und so weiter.

00:08:45 Sprecher 2

Das war noch ganz früh, aber es war eigentlich klar, dass die Informatik
wie eben zum Beispiel Elektrotechnik ein wichtiges und ein praktisches
Gebiet ist.

00:08:56 Sprecher 2

Also da ich glaube, kann mich nicht daran erinnern, dass sich
irgendjemand ein negativen Kommentar gemacht hat.

00:09:06 Sprecher 1

Schon gesagt, es gab viele Sachen zu dieser Zeit noch nicht.

00:09:10 Sprecher 1

Wie konnte man sich das Studium dann an der ETH vorstellen?

00:09:14 Sprecher 1

Also womit habt ihr gearbeitet?

00:09:16 Sprecher 2

Also das erste Mal ist die meiste 75\% der Stunden waren nicht
Informatik, sondern Mathematik und Physik und Elektrotechnik und so
weiter.

00:09:27 Sprecher 2

Ist glaube ich mittlerweile ein bisschen weniger geworden, aber das war
eigentlich die ersten zwei Jahre war mehr so ein Ingenieurstudium mit
Informatik nebendran und dann

00:09:35 Sprecher 2

die zweiten zwei Jahre waren dann effektiv in Informatik.

00:09:40 Sprecher 2

Wir haben angefangen neu in dem Jahr als ich angefangen hatte, war das
neu der Apple 2, der erste Mac Classic mit glaube ich 128 Kilobyte und
einem Floppy Dings.

00:09:59 Sprecher 2

Also es gab so einen großen Raum und da standen paar Dutzend Macs

00:10:05 Sprecher 2

Auf der Seite vom Raum standen noch die Lochmaschinen vom von dem
früheren.

00:10:09 Sprecher 2

Früher hatte man wirklich die Lochkarten gemacht, dann noch im Jahr
vorher, um mit dem Großrechner zu programmieren.

00:10:16 Sprecher 2

Und dort gab es etwas, das war hieß UCSD Pascal und dann haben wir mit
dem Programmieren gelernt.

00:10:24 Sprecher 2

Also Pascal war die die die ETH von der ETH entwickelte Sprache und UCSD
UC San Diego hatte schon ein System geschrieben auf.

00:10:34 Sprecher 2

auf dem auf dem Mac, was so einigermaßen funktionierte.

00:10:39 Sprecher 1

Kam dein Wissen zu dieser Zeit vor allem durch die aus, also durch die
Ausbildung bei der Versicherung und der E.T.A.

00:10:46 Sprecher 1

oder konntest du dich auch selber informieren?

00:10:49 Sprecher 2

Ja, ich war noch einiges in der Bibliothek so nebendran, also es hat mir
ein ein, ich hatte den Vorteil, dass ich schon programmieren konnte,
weil ich das ja anderthalb Jahre oder so gemacht habe und von dem her,
dass die Informatik selber die ersten paar

00:11:04 Sprecher 2

Semester waren easy, weil dort hat man im Programmieren gelernt und so
Basic Algorithms und so weiter und das war, das konnte ich im Prinzip
schon und ich habe einfach auch viel davon Gebrauch gemacht, dass es in
sowohl hier in Liestal in der Bibliothek und wie auch in dann natürlich
in der Lethar Bibliothek verschiedene Bücher gibt, die auch so
verschiedene andere Einführungen hatten, wo man noch mehr Sachen
nebendran lesen konnte.

00:11:32 Sprecher 2

Das das war

00:11:34 Sprecher 2

Wie gesagt, in den ersten 2 Jahren ist nicht so wahnsinnig viel
Informatik.

00:11:38 Sprecher 2

War damals im ETH-Studium, da habe ich auch solche Sachen noch
angeschaut.

00:11:43 Sprecher 1

Gab es dabei eine Lehrperson oder auch einen Studentenkollegen, der dich
dabei besonders geprägt?

00:11:53 Sprecher 2

Hat?

00:11:54 Sprecher 2

Einiges, also das hat an der ETH eigentlich, war ein relativ kleiner

00:12:02 Sprecher 2

Wurst im Prinzip, ich glaube, wir waren am Anfang 180 Leute, aber das
hat sich dann so auf ungefähr 110 oder so relativ schnell
zurückgebildet.

00:12:11 Sprecher 2

Das heißt, man kannte sich noch relativ schnell.

00:12:16 Sprecher 2

Einer der Leute, ja sogar 2, mindestens 2 kenne ich heute noch und und
also einer ist jetzt in Kalifornien, ein anderer ist in in Zürich.

00:12:28 Sprecher 2

Bin ich

00:12:29 Sprecher 2

immer noch in Kontakt und von dem her haben wir da uns da viel
ausgetauscht, auch Sachen zusammen gemacht und aber ich glaube, der
größte, das größte Berufsbild oder Vorbild oder so war natürlich Niklaus
Werth, der war ein Professor, hat auch die Einführungs Teil der
Einführungs Vorlesung selber gemacht und dann die Assistenten, die waren
wirklich sehr gut.

00:12:50 Sprecher 2

Also an der E.T.H., da hast du die Vorlesung, aber da hast du nebendran
noch so nett sich Übungen, aber das

00:12:57 Sprecher 2

Also einfach im Kleineren, das spaltet sich dann auf, all die Leute und
dann hast du einen Assistenten, der da über die verschiedenen Sachen
hergeht und die waren wirklich alle super gut.

00:13:12 Sprecher 1

Perfekt.

00:13:13 Sprecher 1

Also du hattest natürlich einen nach der ETH hattest du mehrere wichtige
Stationen.

00:13:25 Sprecher 1

Meine Mutter meinte sogar, du wärst drei Monate in Norwegen gewesen.

00:13:30 Sprecher 1

Und nach dieser Zeit hattest du auch deinen wichtigen Schritt in die
USA.

00:13:37 Sprecher 1

Wie waren, wie eigneten sich diese Ereignisse?

00:13:41 Sprecher 1

Was war?

00:13:41 Sprecher 2

Ja, also Norwegen war ein Praktikum.

00:13:45 Sprecher 2

Man muss oder musste damals mindestens drei Monate Praktikum machen als
Teil des Studiums.

00:13:52 Sprecher 2

Und ich habe mich für eine

00:13:55 Sprecher 2

Bestimmte Sachen interessiert, objektorientiertes Programmieren war
damals neu und ich habe herausgefunden, dass es das gab es praktisch
nirgendwo in Europa, aber ich habe herausgefunden, dass es in Norwegen
so ein Industrieinstitut gab und an der ETH gab es einige norwegische
Studenten und die haben mir dann geholfen, mit denen Kontakt aufzunehmen
und dann war ich dort eben in einem Praktikum

00:14:24 Sprecher 2

Und das hat mich dann eigentlich bestärkt, dass das Gebiet so was ist,
was wirklich interessant ist.

00:14:30 Sprecher 2

Und gegen Ende vom Studium habe ich mich dann einfach bei verschiedenen
amerikanischen Unis beworben, weil eben in Europa immer noch nicht so
wahnsinnig viel los war in der Gegend, wo ich mich wirklich interessiert
hatte.

00:14:45 Sprecher 2

Also die ETH war eigentlich das Einzige, aber ich wollte nicht unbedingt
an der ETH bleiben, weil das war ein sehr

00:14:53 Sprecher 2

eine Art von von von dem Gebiet und das wurde dann eben sehr geprägt von
der von dem was was die zwei Hauptleute an der ETH machen, war nicht so
meine Sache und von dem her musste ich fast irgendwo anders hingehen im
Prinzip.

00:15:07 Sprecher 2

Und ich setze mich mal da drüben hin.

00:15:15 Sprecher 2

Und das war noch etwas mühsam, weil gab es noch kein Internet, also es
war so ganz ganz am Anfang vom Internet, ich konnte irgendwie einen
E-Mail-Account bekommen.

00:15:26 Sprecher 2

der ETH, aber das waren doch, also die hatten keine E-Mail mehr.

00:15:35 Sprecher 2

Aber das war viel so mit, du schreibst einen Brief und dann zwei, drei
Wochen später schreiben sie was zurück und schicken dir einen Katalog
oder irgend so was Ähnliches und war zum Teil auch widersprüchlich.

00:15:47 Sprecher 2

Also das eine hat gesagt, ja wenn du eben dort doktrinierst, dann
bezahlst du keine Tuition und dann die

00:15:54 Sprecher 2

die Uni hat ihnen was anderes geschickt und gesagt, ja egal was
irgendjemand sagt, du musst zwischen bezahlen und so weiter und so fort.

00:16:01 Sprecher 2

Das war also noch relativ schwierig, also einfach mühsam, das alles
rauszubringen.

00:16:07 Sprecher 1

Es musste dann passieren, dass du den Schritt gewagt hast nach.

00:16:14 Sprecher 2

Weiß nicht, wie du das meinst, also ich meine, also es ist einfach.

00:16:18 Sprecher 1

Aus Unterstützung von der Familie, vielleicht ein Stipendium oder?

00:16:24 Sprecher 2

Also ich hatte dann noch parallel mich beworben für ein Stipendium von
den Amerikanern, das heißt das Fulbright Program.

00:16:34 Sprecher 2

Es gab mal einen Typ, der hieß Fulbright als Nachname und das
unterstützt Austauschstudien.

00:16:42 Sprecher 2

Das habe ich dann auch bekommen.

00:16:44 Sprecher 2

Das hilft, aber das ist unabhängig sozusagen von der Uni.

00:16:47 Sprecher 2

Also das zahlt dir dann die Tuition, wenn nötig.

00:16:55 Sprecher 2

Aber es wäre auch, glaube ich, auch ohne gegangen.

00:16:57 Sprecher 2

Wäre auch schlussendlich einfacher gewesen, weil das hat dann später
alle möglichen Visa-Probleme verursacht.

00:17:07 Sprecher 1

Soweit ich weiß, ging es dann nach Stanford.

00:17:12 Sprecher 1

Was war denn eigentlich dein Plan B, wenn das jetzt nicht geklappt
hätte?

00:17:16 Sprecher 2

Also ich hatte mich bei mehreren Unis beworben.

00:17:21 Sprecher 2

Und mein Plan B war wahrscheinlich UC Davis, das ist dort auch ganz in
der Nähe.

00:17:30 Sprecher 2

Die haben auch einen Professor gehabt, der was Interessantes gemacht.

00:17:37 Sprecher 2

Und ich weiß nicht mehr genau, es war noch was in Oregon oder vielleicht
San Diego, eines von.

00:17:48 Sprecher 1

Denen.

00:17:50 Sprecher 1

Die USA, die

00:17:51 Sprecher 1

Also wie du es schon angedeutet hast, der Briefkontakt dauerte auch
etwas.

00:17:57 Sprecher 1

Was bedeutet das im Kontext zur Familie, dem Kontaktaustausch mit deiner
Familie?

00:18:05 Sprecher 1

Wie habt ihr das für euch geklärt?

00:18:08 Sprecher 2

Das war damals ja eigentlich mühsam.

00:18:11 Sprecher 2

Also weil Telefonanruf kostete 5 Franken pro Minute.

00:18:17 Sprecher 2

Es war also nicht wirklich praktikabel und da hat man einfach

00:18:21 Sprecher 2

haben zum Brief geschrieben oder erhalten und dann kam so langsam E-Mail
gerade so um die also es war ich ging 1988 und der Jörg z.B hatte schon
ein E-Mail Konto hier in der Schweiz irgendwo und in Stanford hat dann
jeder Student eine E-Mail gehabt und das hat dann geholfen weil ich
glaube mein Vater hat sich dann auch irgendwo ein E-Mail Account
geschafft wo er bei der Arbeit

00:18:51 Sprecher 2

reinschauen konnte.

00:18:52 Sprecher 2

Da konnte man also auch so noch Kontakt haben, aber sonst war es im
Prinzip Brief oder ein Besuch.

00:19:01 Sprecher 2

Ich glaube, ich kam erst nach dem ersten Jahr zurück in die Schweiz für
drei Wochen oder so.

00:19:12 Sprecher 1

Dein erstes Ankommen in die USA war das, wie war, es ist

00:19:19 Sprecher 1

Die USA unterscheidet sich ja ziemlich stark zwischen der Schweiz.

00:19:23 Sprecher 1

Du hast schon erwähnt, dass es einen fachlichen Unterschied gibt.

00:19:29 Sprecher 1

Was war so der größte Kulturschock, fachlich und privater?

00:19:35 Sprecher 2

Also ich meine, man muss gewisse Sachen einfach lernen.

00:19:40 Sprecher 2

Also die Uni funktioniert da ganz anders als hier.

00:19:44 Sprecher 2

weil du eben den Campus hast und Undergrad and Grad funktioniert
sozusagen anders.

00:19:53 Sprecher 2

Schon nur bezahlen.

00:19:55 Sprecher 2

Also Konto bekommen, wie das Rechnung bezahlen ist, ganz anders und muss
man halt einfach lernen.

00:20:02 Sprecher 2

Für mich war das, ich würde sagen, die zwei größten Überraschungen waren
die Infrastruktur für Studenten bei Stanford war eigentlich

00:20:13 Sprecher 2

schlechter als in der ETH.

00:20:16 Sprecher 2

Und das heißt wir hatten als neue PhD Studenten nur so ein paar
Terminals, also wirklich so ASCII Terminals, was dann auf einen Unix
Rechner ging und in der Schweiz zu dem Zeitpunkt hattest du dann schon
so Sun Workstations, grafisches Display, also viel viel netter.

00:20:35 Sprecher 2

Und das war eigentlich absichtlich, weil sie wollten, dass

00:20:38 Sprecher 2

dass man sich sozusagen so schnell wie möglich an einen Professor
anschließt, an eine Forschungsgruppe anschließt.

00:20:45 Sprecher 2

Und die haben dann das Fernsehequipment sozusagen, nicht?

00:20:50 Sprecher 2

Also das war also sozusagen der Grund Dings ohne die, dass man in der
Forschung ist, war viel weniger ausgestattet als in der Schweiz, wo das
wirklich breit gut ausgestaltet war.

00:21:01 Sprecher 2

Aber das andere war eigentlich für mich mehr ein kulturelles Ding.

00:21:05 Sprecher 2

Die Amerikaner sagen eigentlich nicht,

00:21:08 Sprecher 2

oder also sie sagen schon, was sie denken, aber man hört es nicht, wenn
man aus der Schweiz kommt.

00:21:13 Sprecher 2

Und das musste ich am Anfang noch lernen, wie wie Feedback funktioniert,
weil das das funktioniert einfach ganz anders.

00:21:26 Sprecher 1

Deine Dissertation drehte sich vor allem um die Programmiersprache Self.

00:21:34 Sprecher 1

Was war da eigentlich dein

00:21:36 Sprecher 1

Gedanke daran erstens es wieso wieso die Programmiersprache self und
jetzt auch noch mal wo du erwähnt hast die Forschungswelt dass die auch
anders ist als in der ETH wie war das dort nochmals wie wie unterschied
sich das?

00:21:59 Sprecher 2

Das Hauptding war eigentlich nicht unbedingt die Sprache sondern wie man
sie implementiert

00:22:06 Sprecher 2

Und die Sprache war interessiert, weil es interessant, weil sie
eigentlich sehr, sehr schwierig ist.

00:22:11 Sprecher 2

Also sie ist sehr, sehr, sehr ein reines Konzept und von dem her, wenn
man sie sozusagen normal implementiert, ist es super langsam.

00:22:18 Sprecher 2

Und ich hab mich, das war eigentlich, dieser Professor war eigentlich
der Grund auch, wieso ich mich bei Stanford beworben hab, weil er hatte
seinen Ph.D.

00:22:27 Sprecher 2

in Berkeley gemacht, war relativ neu in Stanford und er hat damals dort
sehr interessante Arbeit gemacht, eben in dem Bereich, wie macht man
solche Sachen.

00:22:35 Sprecher 2

schnell, obwohl es sehr pur ist.

00:22:38 Sprecher 2

Und das war eigentlich das, was mich interessiert hat.

00:22:40 Sprecher 2

Also mehr Compiler als die Sprache selber.

00:22:42 Sprecher 2

Und der hatte eine kleine Gruppe, waren glaube ich drei andere Studenten
drin und war interessiert daran, jemand neu zu tun, einen neuen zu
haben.

00:22:58 Sprecher 2

Also da war das Stipendium vielleicht in den ersten zwei Jahren
hilfreich, weil er musste aber nichts bezahlen.

00:23:05 Sprecher 2

für mich weil ich vom Stipendium bezahlt wurde und das hat dann so
angefangen also ich habe dann einfach angefangen mich da einzuarbeiten
und du bekommst dann nicht unbedingt einen Auftrag also mindestens dort
war es nicht so dass dann der Dave der Dave war mein mein Professor der
hat nicht gesagt hey wo ist dein Thema für deine dies ist X sondern
einfach ja do something

00:23:35 Sprecher 2

Und am Anfang einfach "Hilf mal hier" und ich hab das meine
rumgeschrieben und alles Mögliche.

00:23:43 Sprecher 1

Fühlt es sich damals .

00:23:46 Sprecher 1

Also, zu dieser Zeit, als du dort warst, fühlte es sich genau richtig an
für dich?

00:23:52 Sprecher 1

Also, was du gemacht hast und auch wo du das machst.

00:23:57 Sprecher 2

Also, ich, Armin .

00:23:59 Sprecher 2

Jain, also Stanford ist ein .

00:24:02 Sprecher 2

ein toller Platz im Sinne, sehr schöner Campus.

00:24:05 Sprecher 2

Kalifornien ist nett, du kannst mit dem Velo überall hingehen.

00:24:11 Sprecher 2

Die Uni hatte wirklich sehr gute Vorlesungen und Leute und so weiter.

00:24:16 Sprecher 2

Also da hast du viele Sachen gelernt, hast auch so Seminare gehabt, wo
du einfach reinschauen kannst und über wirklich ganz neue Sachen.

00:24:27 Sprecher 2

Und die Gruppe, wo ich war,

00:24:31 Sprecher 2

hat mir eigentlich auch gefallen.

00:24:32 Sprecher 2

Das war allgemein Platznot, also ich hatte nicht wirklich ein Büro oder
auch nur einen Desk.

00:24:39 Sprecher 2

Bin da sozusagen schwarz gefahren eine Weile und aber ich habe mein
eigentliches Thema nicht gefunden, also nicht im ersten Jahr gefunden.

00:24:53 Sprecher 2

Das war irgendwie nach anderthalb Jahren, wo so die Hauptideen langsam
rauskommen.

00:25:02 Sprecher 1

Soweit ich weiß, bist du dann nach Santa Barbara gegangen.

00:25:11 Sprecher 1

Ich glaube auch dort hast du deine Frau kennengelernt oder war das schon
zuvor?

00:25:16 Sprecher 2

Das war schon zuvor.

00:25:20 Sprecher 1

Wie hat sie dich geprägt?

00:25:23 Sprecher 2

Also ich habe sie schon am Anfang in Stanford kennengelernt, aber sie
kam

00:25:29 Sprecher 2

unabhängig, habe Biologie studiert und war zum Teil weg, weil sie musste
noch, also sie hat ihre Masters in Hawaii gemacht, also die Arbeit
sozusagen und dann auch noch ein Masters in Deutschland weggemacht.

00:25:47 Sprecher 2

Das war also immer so ein bisschen in Amerika sagt man "two body
problem", das heißt, du bist nicht am selben Ort, ein paar aber nicht am
selben Ort.

00:25:59 Sprecher 2

immer und war damals schon ein bisschen einfacher weil E-Mail war so
einigermaßen funktioniert es außer an den deutschen Unis da war das noch
nicht bekannt auch zwei drei Jahre später und war für mich auch ein eine
wie soll man ein Eintritt sozusagen in eine ganz andere Welt weil das
war waren die Biologiestudenten und war noch eigentlich sehr interessant

00:26:28 Sprecher 2

einfach den Kontrast zu haben auch.

00:26:30 Sprecher 2

Und zum Beispiel am Anfang, die haben alle gelächelt, dass wir so viel
auf E-Mail gemacht haben, weil das war unpraktisch und odd, nicht?

00:26:41 Sprecher 2

Also das macht man einfach nicht, nicht?

00:26:44 Sprecher 2

Und ein Jahr später haben sie dann Macs bekommen.

00:26:47 Sprecher 2

Also jeder hat einen Mac, Desktop Mac und dann war alles auf E-Mail.

00:26:52 Sprecher 2

Bei ihnen auch, nicht?

00:26:54 Sprecher 2

Also da kann man wirklich sehen, wie diese sozusagen

00:26:58 Sprecher 2

diese Arbeitsweise sich verbreitet hat, ziemlich rasant von eben ein
Jahr vorher war es, alle haben gesagt, völlig unpraktikabel, braucht es
nie.

00:27:07 Sprecher 2

Und ein Jahr später läuft das ganze Projekt plötzlich auf E-Mail.

00:27:12 Sprecher 2

Nur weil eben die Macs billig genug geworden sind, dass auch die
Biologen einen haben können.

00:27:18 Sprecher 2

Und wenn dann alle Macs vernetzt sind und dann ist E-Mail eine ganz,
also eine sehr praktische, ein praktisches Tool.

00:27:28 Sprecher 1

In Santa Barbara strebt es ja eigentlich vor allem deine Professur an.

00:27:35 Sprecher 1

Es kam dann aber auch noch zum Kontakt mit Larry Page und Sergey Brin.

00:27:42 Sprecher 1

Wie kam das eigentlich zustande, dass ihr euch so gefunden habt?

00:27:47 Sprecher 2

Also ich war ein Teil davon.

00:27:51 Sprecher 2

Also das Problem war ja, dass ich in Santa Barbara war und dann

00:27:55 Sprecher 2

war immer noch in Stanford und ich habe auch noch in einem Startup also
ich hatte bei meinem Abschluss dann ein noch ein Startup gegründet mit
jemand mit zwei Leuten aus New York und dem Lars den du vielleicht
kennst aus Dänemark und und dann nachher noch zwei drei anderen Leuten
und das war das heißt ich ging immer hin und her und nach

00:28:23 Sprecher 2

fünf Jahren wollte ich, da kam so langsam die Möglichkeit, ein Spadical
zu machen, von also einem Jahr wegzugehen von der Uni und ich wollte was
tun, was eben in der Gegend von Paralto ist, weil damit ich da nicht
jedes Wochenende hin und her fliege.

00:28:45 Sprecher 2

Das wurde langsam mühsam.

00:28:47 Sprecher 2

Und damals war der Internetboom

00:28:54 Sprecher 2

äh also da komm war wegen vollem Gang und du bekommst ständig eine
E-Mail oder oder einen Anruf oder sowas weißt du äh von irgendwelchen
recruiters so im Stil arbeite für x.com und bei.com und also alles
konnte ich hier nicht und die meisten Sachen waren aber nicht also
einfach nicht interessant oder weil die das war irgend so ein

00:29:19 Sprecher 2

E-Commerce Ding oder sowas also einfach nichts technisches zu tun und
dann habe ich zuerst ein ich glaube zuerst war es ein ehemaliger Student
UC Sainabababer Student hat mir davon geschrieben und gesagt ja das
sieht noch interessant aus und dann habe ich aber eigentlich nichts
damit gemacht weil ich verstehe nichts von Suche das

00:29:45 Sprecher 2

Das war damals sehr esoterisch, nannte sich Information Retrieval,
völlig anderes Feld, als ich gemacht habe, weil ich machte so Compiler,
Computer Architecture, Performance, solche Sachen.

00:29:57 Sprecher 2

Und dann nachher hat noch, ich war jeweils einen Tag bei Stanford.

00:30:02 Sprecher 2

Ich hatte ein Projekt mit einer Professorin in Stanford.

00:30:06 Sprecher 2

Und die hat dann auch gesagt, no, no, you should talk to Lions regular,
die haben Probleme für jemand wie dich.

00:30:15 Sprecher 2

Und dann habe ich ihnen einfach eine E-Mail geschrieben, da hat sich
herausgestellt, die haben mich gekannt, weil ich mal einen Vortrag
gegeben habe, weil es einfach in der in der Audience und you know, one
thing leads another, nicht.

00:30:33 Sprecher 1

Was hast du dann darin, was hast du dann im Startup gesehen, was
vielleicht jemand anderes nicht sah oder

00:30:43 Sprecher 1

Was hat dich denn gefragt?

00:30:44 Sprecher 2

Also im Prinzip, die was ganz anders war, die haben nicht von Geld
geredet und vom vom.

00:30:53 Sprecher 2

Also eben, es ging nicht darum, sondern es ging darum.

00:30:58 Sprecher 2

Also sagen wir mal, also back to back up, nicht?

00:31:02 Sprecher 2

Also zwei Sachen.

00:31:02 Sprecher 2

Das erste Mal, ich hatte mich damals schon daran gewöhnt, dass man

00:31:09 Sprecher 2

um im Internet etwas zu finden fast Informatiker sein muss weil du
musstest die Queries wirklich sehr genau formulieren mit and und Ohr und
Klammern und so weiter damit es nur auch nur ein einigermaßen
funktionierte also im Allgemeinen konntest du Sachen einfach nicht gut
finden nicht und wir haben einfach gedacht okay zum Glück bin ich
Informatiker für mich funktioniert aber für Otto Normalverbraucher wird
es nie funktionieren das ist einfach zu schwierig und ihre Mission

00:31:39 Sprecher 2

das können wir ändern.

00:31:41 Sprecher 2

Also hier ist all diese Informationen auf dem Internet und niemand kann
sie finden und wir können das ändern.

00:31:49 Sprecher 2

Und schon damals auch das Demoprojekt sozusagen war wirklich viel viel
besser als was es sonst gab.

00:31:58 Sprecher 2

Also Beispiel, du konntest Stanford eintippen und das erste Resultat war
die Stanford Homepage.

00:32:04 Sprecher 2

Und bei allen anderen Suchmaschinen kam alles mögliche, aber nicht die
Stanford Homepage.

00:32:09 Sprecher 2

Das kann man sich heute nicht vorstellen.

00:32:11 Sprecher 2

Aber das war wirklich ein Tag und Nachtunterschied.

00:32:13 Sprecher 2

Und dann das Zweite war, dass all die Leute, die ich getroffen habe,
also Larry Sergej und die ersten zwei, drei Angestellten, die waren
einfach sehr gut.

00:32:22 Sprecher 2

Also wirklich technisch gut.

00:32:25 Sprecher 2

Und persönlich auch nett.

00:32:30 Sprecher 2

Und es war wirklich klar, dass man dort etwas lernen könnte.

00:32:34 Sprecher 2

Und mein .

00:32:37 Sprecher 2

Mein also war nicht unbedingt geplant, aber meine Annahme war.

00:32:40 Sprecher 2

okay ich mache das jetzt ein Jahr während ich da auf die GISG warte und
nachher entscheiden wir dann wohin wir gehen und ich lerne sicher was
weil das ist was Neues und ist interessant habe damals nicht gewusst was
war noch nie in einem Datenzentrum und beim zweiten Vorstellungsgespräch
nahm mich der Larry mit in das lokale Datenzentrum wo wir was gemietet
haben und sah interessant aus so why not nicht.

00:33:07 Sprecher 1

Also du hast das

00:33:09 Sprecher 1

allgemeine Risiko auch relativ niedrig eingeschätzt, da es auch für dich
gerade perfekt reingepasst hätte.

00:33:16 Sprecher 2

Also ich hatte, das war nicht, sagen wir es mal so, ich hab das mehr so
als Lehrjahr, als Lehrjahr angeschaut, nicht die, dass dass der Lohn war
sehr niedrig.

00:33:32 Sprecher 2

Also ich konnte davon meine Steuern bezahlen, mehr oder weniger, aber
ich hatte ein gutes Einkommen von

00:33:39 Sprecher 2

Also der Startup den ich da erwähnt hatte, der wurde von Sun
Microsystems gekauft und ich war dort noch Consultant und hatte auch ein
Einkommen von dort und dann das Uni-Einkommen.

00:33:54 Sprecher 2

Du wirst während des Badaco trotzdem noch bezahlt.

00:33:57 Sprecher 2

Also sozusagen alle sechs Jahre verdienst du dir ein Jahr oder drei,
drei, neun Monate.

00:34:02 Sprecher 2

Von dem her musste ich nicht davon leben.

00:34:07 Sprecher 2

Und das war nicht riskant im Sinne von: Ich habe nichts davon erwartet.

00:34:15 Sprecher 2

Und ich musste nichts davon erwarten.

00:34:16 Sprecher 2

Also ich musste nicht von dem leben.

00:34:21 Sprecher 1

Folgendermaßen wurde wahrscheinlich der Schritt zu Hause gut
aufgenommen.

00:34:27 Sprecher 2

Ich glaube nicht, dass das so klar verstanden wurde unbedingt, weil
Google hat niemand gekannt.

00:34:35 Sprecher 2

Es hat auch auf Deutsch noch nicht wirklich funktioniert.

00:34:38 Sprecher 2

Das kam so erst etwa ein Jahr später, weil der Index, also welchen Teil
vom Web Google effektiv gekannt hat, war es recht klein und vor allem
USA dominiert.

00:34:50 Sprecher 2

Das heißt, auf Englisch hat es gut funktioniert.

00:34:52 Sprecher 2

Auf Deutsch hat es nicht gut funktioniert.

00:34:55 Sprecher 2

Erst einmal, weil wir ein paar Fehler gemacht haben, weil auf Deutsch
ist anders als Englisch, aber vor allem, weil es einfach sehr wenig
deutsche

00:35:04 Sprecher 2

Seiten gehabt habe nicht und da kannst natürlich du kannst nicht was
finden wo von dem du nichts weißt oder also das Ranking kann das nicht
fixen das heißt Google wurde glaube ich in Deutschland erst so 2000
ungefähr also anfangs 2000 langsam bekannt und von dem her glaube ich
haben sie ich war ja vorher schon an dem anderen Startup und war

00:35:32 Sprecher 2

ja auch immer noch Professor.

00:35:33 Sprecher 2

Also ich weiß nicht, wie viel das wirklich eine Frage war.

00:35:36 Sprecher 2

Und ein paar paar Jahre später, als dann Google auf auf Deutsch
funktioniert, dann war es, dann haben alle gefunden, oh, das ist ja
toll, weil es funktioniert einfach viel besser, nicht?

00:35:45 Sprecher 2

Also der Consumer Value war wirklich da, nicht?

00:35:48 Sprecher 2

Das haben wir dann schon gesehen.

00:35:51 Sprecher 1

In den Anfangsjahren gab es wahrscheinlich paar Ausfälle und auch paar
Momente, in denen man, ich sag mal, Hotfixes improvisieren musste.

00:36:04 Sprecher 1

Gab es jetzt vielleicht einen Moment, an den du dich wirklich innerst
und der dich auch geprägt hat von seiner Situation her vielleicht.

00:36:12 Sprecher 2

Ich glaube nicht einen einzigen, weil da war wirklich ständig alles
kaputt im Prinzip.

00:36:18 Sprecher 2

Also die, ich würde sagen, wir hatten am Anfang mindestens einmal pro
Woche eine, wahrscheinlich mehrmals pro Woche.

00:36:28 Sprecher 2

Das war also fast ein Dauerzustand, weil

00:36:32 Sprecher 2

Das kam aus der Uni raus und das war sozusagen eine gute Idee, aber
nicht wirklich professionell geschrieben und auch nicht konzipiert für
auch schon nur eine Million Benutzer zu haben.

00:36:48 Sprecher 2

Das war ja ein Forschungsprojekt, da hat man nicht eine Million
Benutzer.

00:36:52 Sprecher 2

Und von dem her praktisch jede Woche war wirklich die Frage, wie
überleben denn wir den nächsten Montag?

00:36:59 Sprecher 2

Montag war das höchste

00:37:02 Sprecher 2

Traffic der Woche und damals typischerweise jeder Montag war 5\% mehr
als der Montag vorne da war Montag übrigens weil die Leute übrigens
Wochenende hatten sie neue Fragen aber zu Hause hatten sie vielleicht
noch kein Internet oder nur Modem oder so und dann sind sie am Montag in
die Schule oder an Work gekommen und dort gab es besseres Internet und
haben sie dann dort die die Frage gemacht und z.B typisch ich würde
sagen 1999

00:37:32 Sprecher 2

sicher im ersten halben Jahr am Montag um 12 Uhr, hat es vielleicht
funktioniert.

00:37:42 Sprecher 2

Wenn es funktioniert, hat es fünf Sekunden genau meine Antwort bekommen
und so weiter.

00:37:48 Sprecher 2

Das war also wirklich eigentlich mein erster, man konnte sich einen
Jobtitel selber suchen, mein erster Jobtitel war "Search Engine
Mechanic" und

00:37:59 Sprecher 2

Ich wollte nicht, dass die Leute fragen, "Why is a search engine
mechanic?" Meine Antwort war "because everything is broken".

00:38:06 Sprecher 2

Und das hat auch gestimmt.

00:38:08 Sprecher 2

Also das war wirklich sehr .

00:38:12 Sprecher 2

Da fliegt einfach alles nicht.

00:38:16 Sprecher 1

In den Anfangs 2000er gab es ja einen großen Moment, zum Beispiel die
Dotcom-Krise.

00:38:28 Sprecher 1

Wie war

00:38:29 Sprecher 1

zu dieser Zeit die Stimmung auch bei dir gab es also warst du schon
zuversichtlich dass Google da eher als Gewinner rauskommt oder war das
noch alles sehr ungewiss?

00:38:39 Sprecher 2

Ja also glücklicherweise wurde Google gerade zum richtigen Zeitpunkt
gegründet sozusagen früh genug dass man noch Geld also Mitte 99 hatten
sie Funding Round

00:38:54 Sprecher 2

dass das noch funktionierte, aber spät genug, dass man schon gesehen
hat, dass die Bubble wirklich eine Bubble ist.

00:38:59 Sprecher 2

Also das war nicht unbedingt eine Überraschung.

00:39:02 Sprecher 2

Und wir haben eigentlich von Anfang an Geld gespart, weil also das war
relativ diszipliniert, weil wir nicht so wahnsinnig viel Funding hatten
und weil eben das Problem ein Problem ist, wofür man viele Rechner
braucht, nicht?

00:39:17 Sprecher 2

Und dann, es war nicht klar, dass wir das Geld bekommen, um wirklich so
viele Rechner zu haben.

00:39:22 Sprecher 2

Und

00:39:25 Sprecher 2

Da kommt Blase ist geplatzt im Mai/Juni 2000.

00:39:31 Sprecher 2

Das war schon sehr dramatisch und wir hatten aber noch, wir waren noch
nicht profitabel, aber wir hatten so einigermaßen eine Idee, wie man das
innerhalb von einem bis zwei Jahren machen könnte.

00:39:48 Sprecher 2

Aber das war wirklich dramatisch.

00:39:50 Sprecher 2

Wir waren damals in Mountbillo in unserem ersten

00:39:52 Sprecher 2

Büro und innerhalb von zwei und das war so ein Office Park, also gross
viele Gebäude, dutzende von Gebäuden, ziemlich busy, also jeden Morgen
Stau oder so ähnliches.

00:40:04 Sprecher 2

Und acht Wochen später war es eine Wüste.

00:40:09 Sprecher 2

Also wirklich jedes du fährst durch und nichts ist auf dem Parkplatz.

00:40:16 Sprecher 2

Und es gab ein Haus, also ein Büro, wo es noch

00:40:21 Sprecher 2

noch Leute gehabt, das war ein eben nicht eine Dotcom Firma, sondern ein
ein Mutual Fund, also eine Finanzfirma.

00:40:30 Sprecher 2

Und sonst war einfach alles weg.

00:40:32 Sprecher 2

Es war wirklich wie wenn eine Bombe eingeschlagen hat und alle Leute
sind ausgezogen innerhalb von Wochen.

00:40:39 Sprecher 2

Es war wirklich dramatisch.

00:40:40 Sprecher 2

Aber wir waren, das war für uns eigentlich gut, weil da haben wir
plötzlich viele gute Leute gehabt, die haben einen Job besucht und es
war dann viel einfacher.

00:40:50 Sprecher 2

sehr gute Leute anzustrengen als vorher.

00:40:55 Sprecher 1

War dir vorher vielleicht schon bewusst, oder genau halt eben in diesem
Moment, dass Google doch das Potenzial hat, so groß zu werden.

00:41:05 Sprecher 2

Nein, es ging eigentlich mehr ums Überleben.

00:41:08 Sprecher 2

Das Problem war, es war ganz klar, dass die Funktionalität überleben
wird.

00:41:14 Sprecher 2

Es war dann zu dem Zeitpunkt, war es wirklich,

00:41:19 Sprecher 2

die einzige Art etwas zu finden praktisch auf der Mitte.

00:41:23 Sprecher 2

Es gab einen einen Competitor, über den wir uns Sorgen gemacht haben.

00:41:26 Sprecher 2

Die waren in Norwegen, aber die wurden dann auch abgehängt, nicht?

00:41:30 Sprecher 2

Und von dem her ist es etwas was nützlich ist und wird das irgendwie von
jemandem gekauft vielleicht oder sowas von EOL oder sowas ähnliches.

00:41:42 Sprecher 2

Also wenn wenn wir nicht selber überleben, das war eigentlich das war
eigentlich klar, nicht?

00:41:48 Sprecher 2

Wird es

00:41:49 Sprecher 2

Kommerziell etwas Großes war überhaupt nicht klar, weil wir hatten nicht
eine, also auf Englisch sagen wir mal monetize, weiß nicht wie das auf
Deutsch heißt.

00:41:59 Sprecher 2

Also wir hatten einfach noch nicht das richtige Ad-Konzept gefunden für
die Suche und deswegen waren die Einnahmen vor allem Lizenzeinnahmen und
so weiter und ein paar Ad-Revenue, aber das war einfach nicht

00:42:18 Sprecher 2

das hat knapp gereicht um nicht bankrott zu gehen sozusagen und das kam
dann glaube ich ungefähr 2003 also das war die dritte Version von
unserem Ad System was dann wirklich das richtige Rezept gefunden hat und
dann ging es eigentlich schlagartig nach oben.

00:42:37 Sprecher 2

Also dann war es klar, dass wir profitabel sind.

00:42:43 Sprecher 1

Jetzt auch du hast geredet wie du

00:42:48 Sprecher 1

jetzt auch selber genannt hast als jemanden, der eigentlich vor allem
die Box fixt, also die Engineer und das und das auch vor allem darum
geht, dass um die Optimierung der Suche.

00:43:05 Sprecher 1

Heute läuft ja jetzt ein Großteil der der Suche im Internet über eure
Infrastruktur und über eure Software.

00:43:20 Sprecher 1

Wie fühlt es sich jetzt momentan eigentlich an?

00:43:23 Sprecher 1

Und wann wurde dir eigentlich bewusst, wie groß dieses Projekt ist?

00:43:29 Sprecher 2

Ja, das kam vielleicht.

00:43:32 Sprecher 2

Es ist ein bisschen schwierig, weil es hat einfach nicht aufgehört.

00:43:35 Sprecher 2

Also ich glaube, als Mensch kann man sich nicht gut erlauben.

00:43:39 Sprecher 2

Ich habe eigentlich immer schon gesagt, das müsste irgendwann .

00:43:43 Sprecher 2

Also immer schon, sagen wir mal 2002, 2003 haben sie gesagt, das müsste
eigentlich eine 100 Milliarden Firma sein.

00:43:52 Sprecher 2

Und aber ich habe das eigentlich nicht, weil das geht, das ging ja 10,
20 Jahre, nicht?

00:44:02 Sprecher 2

Und ich glaube, man kann sich einfach schlecht vorstellen, was Zins und
Zinseszins sozusagen in 20 Jahren macht, nicht?

00:44:11 Sprecher 2

Also wenn du jedes Jahr

00:44:12 Sprecher 2

um 30\% wächst, dann nach 20 Jahren ist es ein riesen Wachstum.

00:44:18 Sprecher 2

Aber es fühlt sich nicht an, dass du irgendwie einen Sprung machst,
nicht von A nach B und vorher warst du klein, jetzt bist du plötzlich
groß, sondern das ist wirklich graduell.

00:44:27 Sprecher 2

Und von dem her war es vielleicht erst in 2010er Jahren, wo es dann
wirklich klar war, dass es eine Großfirma wird und nicht einfach eine
kleine Firma, die wächst.

00:44:41 Sprecher 2

Und es gab ja viele Sachen auch noch nicht.

00:44:43 Sprecher 2

Also zum Beispiel YouTube war erst so ungefähr, sagen wir mal 2008/9
wirklich was, was funktionierte.

00:44:51 Sprecher 2

Android und iPhone kam auch gerade um die Zeit.

00:44:54 Sprecher 2

Also das ganze Internet war völlig anders überhaupt.

00:44:56 Sprecher 2

Mobil, eben nichts mobil vorher.

00:44:59 Sprecher 2

Also hat sich immer entwickelt und dann eben auch den Dings verstärkt.

00:45:07 Sprecher 1

Du gingst ja eigentlich von der Forschung

00:45:11 Sprecher 1

mehr dann zu zum Führen vom Mitarbeiter eigentlich über auch der Praxis
ja du hattest mit dem Wachstum von Google auch immer wie mehr
Verantwortlichkeit über über mehr Mitarbeiter bist du hast du das eher
auf dem Weg versucht Inge mitzulernen

00:45:38 Sprecher 1

Oder gab es eine Person vielleicht in deinem Kreis, mit dem du dich die
ganze Zeit abgesprochen hast?

00:45:44 Sprecher 1

Wie hast du das Fieren von Personen gelernt?

00:45:47 Sprecher 1

Und gab es vielleicht einen Punkt, wo du oder einen Themenbereich, den
du besonders schwer fandest?

00:45:55 Sprecher 2

Also das war relativ schwierig.

00:45:59 Sprecher 2

Ich wurde ja zum ersten VP im Sommer von 1999.

00:46:03 Sprecher 2

Das war im Wesentlichen, weil ich der Einzige war, der

00:46:06 Sprecher 2

auch irgendwo ein bisschen Managementerfahrung hatte als Professor im
Prinzip und vom Startup und ich habe gesagt okay mache ich aber nur wenn
wir einen wirklichen VP anstellen der das der effektiv weiß wie man das
macht und das hat das dauert ja dann ungefähr 18 Monate weil am Anfang
war es schwierig jemanden zu finden weil da kommen wir immer noch also
alle guten Leute hatten massenhaft Job Office und so weiter und nach 18
Monaten hatte ich dann so

00:46:35 Sprecher 2

glaube ich Engineering war vielleicht 150 Leute und das äh da war ich
froh das abzugeben das hat also nicht unbedingt so viel Spaß gemacht
weil ich einfach noch nicht wirklich gelernt hatte wie wie wie man das
besser macht nicht und der der erste wirkliche VP der war den habe ich
sogar noch gekannt von von von den 90er Jahren von Sunlabs

00:47:04 Sprecher 2

Und der hatte, der war das Gegenteil von mir, der war über 60 und hatte
das schon seit 35 Jahren gemacht.

00:47:13 Sprecher 2

Von dem her habe ich viel dann von dem gelernt.

00:47:16 Sprecher 2

Einfach du arbeitest zusammen, schaust ihm zu und man kannst eben auch
diskutieren.

00:47:20 Sprecher 2

Wir haben dann in den nächsten zwei Jahren auch nochmals zwei, drei
Leute gefunden, die auch viel Erfahrung hatten, viel mehr als ich.

00:47:33 Sprecher 2

Also Wayne Rosing war eben dieser erste, Bill Korn kam etwas später und
die von denen habe ich doch viel gelernt.

00:47:41 Sprecher 2

Und dann also als der Wayne Anfang ging von 170 auf oder weiß ich was es
war auf 15 und dann dauerte es zwei drei Jahre bis ich wirklich bis ich
wieder zurück war auf so ungefähr was ähnliches und das hat sich ganz
anders angefühlt.

00:47:59 Sprecher 2

Weil ich hatte eben Zeit was zu lernen und

00:48:03 Sprecher 2

Weißt du, wenn du irgendwas machst und bist nicht gut, dann ist es mehr
so stressig und overwhelming oder weiß ich was und braucht sehr viel
Zeit, weil du eben einfach im Prinzip ein Anfänger bist und dann die 3
Jahre später war es hat sich ganz anders hat sich viel komfortabler
angefühlt und war war auch nicht unbedingt eine riesige Belastung im
also dass alle Zeit nimmt und dass man dann keine Zeit mehr hat für
technische Sachen.

00:48:35 Sprecher 1

So ein großes Thema in Google war eigentlich auch noch, dass ihr 2007
angekündigt habt, ihr werdet klimaneutral.

00:48:48 Sprecher 1

Woher kommt eigentlich dein persönliches Interesse an dieser
Nachhaltigkeit und es gibt ja jetzt wieder das Ziel auf 2030
klimaneutral zu werden.

00:48:59 Sprecher 1

Wie ist das so für dich?

00:49:01 Sprecher 2

Also das sind zwei verschiedene Begriffe im Prinzip, weil klimaneutral
kann ja verschiedene Sachen ändern.

00:49:08 Sprecher 2

Damals waren es eigentlich zwei verschiedene Sachen.

00:49:12 Sprecher 2

Ich glaube schon von meiner Mutter her war, also sie war so für
Naturschutz und so weiter und so weiter.

00:49:23 Sprecher 2

Ich habe ja für den WWF als Kind, da war noch in der Schule und dann hat
man Briefmarken verkauft oder weiß ich fast nicht und eben um

00:49:31 Sprecher 2

Flamingos oder oder oder Pandas oder was auch immer, Gorillos zu
schützen.

00:49:36 Sprecher 2

Und dann Giske hat ja Ökologie studiert und eben auch die Leute, die ich
dann in dem Kreis getroffen habe.

00:49:44 Sprecher 2

Also von dem her hatte ich schon eine eine, wie sagt man, eine
Affinität, eine eine Beziehung zu das zu dem und dann Langenzeuge zum
Beispiel waren auch also das war eigentlich damals schon klar, dass das
Klimaproblem

00:50:00 Sprecher 2

eines der größten Probleme sein wird.

00:50:02 Sprecher 2

Und von dem her war das nicht irgendwie, wie ich dir gesagt habe, außer
ich war der einzige, der sagte: "Hey, wir sollen was tun", sondern das
war eigentlich wirklich ein Haufen Konsensus.

00:50:11 Sprecher 2

Und das hat sich eigentlich besser entwickelt, als ich gedacht habe.

00:50:17 Sprecher 2

Wir haben 2010, haben wir gesagt, okay, wir möchten so viel erneuerbare
Energie einkaufen jedes Jahr und zwar neue Energien, also neue

00:50:29 Sprecher 2

Solar oder Windparks, wie wir selber verbrauchen.

00:50:35 Sprecher 2

Und wir haben eigentlich gedacht, schaffen wir sicher nicht bevor 2020,
vielleicht sogar noch später.

00:50:41 Sprecher 2

Und dann aber, weil sich Solar und Wind so gut entwickelt hat, war dann
schon in 2017 waren wir dran.

00:50:48 Sprecher 2

Und das war aber nur ein annual, wie man das deutsch sagt, ein annual
matching.

00:50:55 Sprecher 2

Das heißt, wenn wir jedes Jahr x mit

00:50:59 Sprecher 2

Megawattstunden verbraucht haben, dann haben wir auch x Megawattstunden
produziert, aber nicht unbedingt zur selben Zeit und nicht unbedingt am
selben Ort.

00:51:11 Sprecher 2

Das heißt, es ist nicht richtig 100\% renoval und es ist auch nicht
100\% carbon free deswegen.

00:51:18 Sprecher 2

Und dann haben wir den, als wir das erreicht haben, haben wir das
gesagt, okay, wir können eigentlich messen was carbon free heißt und wir
wollen jetzt 100\% annual,

00:51:29 Sprecher 2

auf 100\% 24/7.

00:51:31 Sprecher 2

Das heißt, jedes Datenzentrum oder Büro muss während jeder Stunde von
jedem Tag dieses Matching haben, dass wir also lokal genügend Renoval
Energy haben, das um den Verbrauch abzudecken.

00:51:46 Sprecher 2

Und da sind wir momentan bei etwa 65 bis 70\%.

00:51:51 Sprecher 2

Wir haben schon ein Dutzend Datenzentren, die sind so bei 90\% oder
mehr, aber je nach Region

00:51:59 Sprecher 2

dauert das also noch eine Weile und das ist eigentlich das der Hauptweg
wie wir zu wirklich carbon neutral kommen wollen weil wenn wir das
schaffen für diese Datenzentren dann kann man wirklich sagen diese
Datenzentren verbrauchen für den Strom keinerlei also sind klimatrend
für den Betrieb weil eben weil du irgendwo in der Nähe mit dem Stromnetz
verbunden die

00:52:27 Sprecher 2

24 Stunden pro Tag, die diese Matching Energy heißt.

00:52:32 Sprecher 1

Für diese Nachhaltigkeitsziele, gab es da intern direkt Anklang oder
braucht es dafür eigentlich irgendeine Initiative dafür?

00:52:41 Sprecher 2

Also du musst es immer organisieren, wenn es eine große Firma ist, da
passiert es nicht einfach zufällig.

00:52:46 Sprecher 2

Aber ich glaube eigentlich nicht, dass es groß Kritik daran gegeben hat.

00:52:51 Sprecher 2

Ich glaube, die größte

00:52:57 Sprecher 2

die größte Frage war können wir uns also ist es praktikabel kann man
sich das leisten und und ist es machbar nicht und gewisse Sachen sind da
immer noch nicht ganz 100\% sicher also gerade in Asien zum Beispiel
vulnerable energy in Korea Japan Taiwan ist weit hinten rein und

00:53:26 Sprecher 2

da sind wir momentan am schlechtesten, weil einfach die Umgebung, also
das Land selber hat wenig Möglichkeiten um um und zum Teil auch wenig
politischen Willen um um Renoval Energy zu machen.

00:53:42 Sprecher 1

Du warst ja maßgeblich eigentlich daran beteiligt für den Standort in
Zürich, also was war eigentlich deine Rolle daran

00:53:54 Sprecher 1

Und wieso habt ihr die Schweiz gewählt?

00:53:59 Sprecher 2

Also das war nicht.

00:54:01 Sprecher 2

Also ich bin da mehr so ein Sinnbild als wirklich das, sondern was
passiert ist, wir wollten in Europa ein Entwicklungszentrum aufmachen.

00:54:10 Sprecher 2

Das war so ungefähr 2004.

00:54:11 Sprecher 2

Ich habe ja schon gesagt, unsere Suche auf Deutsch war nicht so gut wie
unsere Suche in den USA, weil nicht, wenn da nur Amerikaner die Sache
entwickeln, dann

00:54:23 Sprecher 2

nächsten nicht mal, was genau.

00:54:30 Sprecher 2

Ja, wir können aufhören, Ruhe zu nehmen.

00:54:34 Sprecher 2

Wir haben noch eine Pause.

00:54:38 Sprecher 1

2000.

00:54:39 Sprecher 2

Also genau, Theorie.

00:54:41 Sprecher 2

Und zwar, da wollten wir eben eine europäische Entwicklungszentrale
haben.

00:54:47 Sprecher 2

Damals waren wir in Dublin schon, aber das war nur so der Betrieb, also

00:54:53 Sprecher 2

Kundenservice und so weiter.

00:54:56 Sprecher 2

Und dann hatten wir 4 Ingenieure, 2 Deutsche, Österreicher und Schwede,
die nach in Kalifornien waren, aber nach Europa zurück wollten.

00:55:08 Sprecher 2

Und die gingen, die haben wir dann auf Tour geschickt, um verschiedene
Orte anzuschauen.

00:55:14 Sprecher 2

Und Zürich war eine der Varianten München, London, weiß nicht, was sonst
noch auf der Liste war.

00:55:23 Sprecher 2

die haben dann empfohlen dass man nach Zürich geht und aus verschiedenen
Gründen waren das die die erste Wahl dort und ich habe meine Rolle war
das helfen zu organisieren weil keiner von denen war ein Manager das
waren einfach Ingenieure und da braucht es halt jemand der da hilft die
Sache effektiv in Gang zu bringen und und und das Büro sozusagen

00:55:53 Sprecher 2

erfolgreich zu machen.

00:55:54 Sprecher 2

Das war dann meine Rolle.

00:56:08 Sprecher 1

2023 bist du dann eigentlich eher vom Management, also um diese Zeit
bist du dann eigentlich eher vom Management zurückgetreten und hast den
Titel Google Fellow angenommen.

00:56:23 Sprecher 1

Gab es dafür einen

00:56:26 Sprecher 1

einen triftigen Grund oder war es einfach oder gab es da einfach eine
Konstellation von Ereignissen, die das ausgelöst haben?

00:56:37 Sprecher 2

Das war so mehr, also das war eigentlich ging über mehrere Jahre, wo ich
Sachen abgegeben habe, weil ich eben das Ziel hatte, auf Null zu kommen
und weil ich habe 20 Jahre, mehr als 20 Jahre gemanagt und das wollte
ich wieder mal

00:56:56 Sprecher 2

was wirklich nur Technisches machen.

00:56:59 Sprecher 2

Und dann in dem Jahr hat es sich so ergeben, dass ich die letzte Person
gefunden habe, wo man das wirklich aufteilen, abgeben kann, ohne dass
ich da gebraucht wurde.

00:57:10 Sprecher 2

Das war ein bisschen Zufall, dass es 2023 war, aber kein Zufall, dass es
passierte.

00:57:17 Sprecher 2

Und dann bin ich dort von einer sehr großen eben Management einer sehr
großen Organisation,

00:57:25 Sprecher 2

Es war eigentlich eben gut Cloud und die Technikinfrastruktur und so
weiter auf, glaube ich, am Anfang auf zwei und mittlerweile auf null
zurückgegangen.

00:57:40 Sprecher 1

Hast du, was sind, du hattest viele Lebensstationen an ganz vielen
Orten.

00:57:48 Sprecher 1

Was ist eigentlich deine Beziehung momentan zu der Schweiz und zu Basel?

00:57:52 Sprecher 1

Also

00:57:53 Sprecher 1

Oder hast du auch einen Punkt, den du als zu Hause definieren würdest?

00:57:58 Sprecher 2

Ja, ist ein bisschen schwierig.

00:58:00 Sprecher 2

Also ich glaube, man liest alles natürlich, wo ich herkomme und was
immer noch eigentlich ein nicht wirklich der Herkunfts oder Heimatort
sozusagen ist.

00:58:14 Sprecher 2

Und ich bin dann, ja, 30 Jahre waren wir in Kalifornien.

00:58:19 Sprecher 2

Das fühlt sich natürlich auch wie eine Heimat an, aber

00:58:23 Sprecher 2

sind ja jetzt auch seit fünf Jahren in Neuseeland.

00:58:26 Sprecher 2

Wenn ich dorthin gehe, fühlt sich das auch sehr vertraut an.

00:58:30 Sprecher 2

Also das ist ein bisschen schwierig zu sagen, weil, also oder zwischen
denen sozusagen zu wählen, nicht?

00:58:38 Sprecher 2

Aber es wäre wahrscheinlich Kalifornien oder Neuseeland, weil ich habe
eigentlich relativ wenig Zeit in der Schweiz,

00:58:51 Sprecher 2

letztes Jahr war es das erste Mal, dass ich längere Zeit hier war und
von dem her lerne ich die Schweiz eigentlich wieder kennen sozusagen.

00:58:58 Sprecher 2

Aber das gab es sicher nicht meine .

00:59:02 Sprecher 2

mein Home.

00:59:04 Sprecher 1

Was wäre etwas, was man aus öffentlichen Quellen über dich nie erfahren
würde oder falsch erfahren würde vielleicht?

00:59:16 Sprecher 1

Also gibt es vielleicht einen Punkt über

00:59:18 Sprecher 1

deine Persönlichkeit oder deine.

00:59:20 Sprecher 2

Vielleicht haben wir einige.

00:59:22 Sprecher 2

Ich habe ich normal, also ich habe mich eigentlich eher rausgehalten,
insbesondere für so persönliche Stories und so weiter.

00:59:29 Sprecher 2

Also man kann es ja nicht vermeiden, wenn du bei meinem Job, bei Google
und so weiter, dass da in der Presse und so weiter kommt.

00:59:36 Sprecher 2

Aber ich habe eigentlich mein persönliches Leben vor allem aus dem
Medienhaus rausgehalten.

00:59:44 Sprecher 2

Also da gibt es eigentlich sehr viele Sachen, die da nicht öffentlich

00:59:47 Sprecher 2

sind also all die privaten Dinge.

00:59:51 Sprecher 2

Ich glaube, dass ich mit Gisky verheiratet bin, ist mit dem Internet
herausfinderbar.

00:59:59 Sprecher 2

Ich glaube schon nur, dass wir keine Kinder haben, das glaube ich nicht,
auch im Internet.

01:00:02 Sprecher 2

Das muss auch gut sein.

01:00:05 Sprecher 1

Hättest du vielleicht eine Geschichte oder eine Anekdote, die du
vielleicht, wenn man an einem Treff ist oder so, über dich erzählen
würdest?

01:00:17 Sprecher 1

um mehr von dir zu hören.

01:00:21 Sprecher 2

Also wie meinst du das?

01:00:23 Sprecher 1

Einfach ein Zelt vielleicht.

01:00:24 Sprecher 2

Oder jemand anderes.

01:00:25 Sprecher 1

Nee, also vielleicht eine Geschichte, die du besonders interessant
findest über vielleicht deine Kindheit, aber auch über deine Zeit im
Studium.

01:00:44 Sprecher 2

Also ich glaube, es war eigentlich

01:00:48 Sprecher 2

Für mich, die Welt war eigentlich, also bis sogar zur ETH nicht, also
damals war die Welt einfach ganz anders, nicht?

01:01:00 Sprecher 2

Und ich habe im Prinzip mein ganzes Leben in Liestal plus minus 10
Kilometer verbracht, so einem typischen Tag, nicht?

01:01:10 Sprecher 2

Bis ich an die ETH ging und schon das fühlte sich eigentlich wie ein
großer Schritt an.

01:01:15 Sprecher 2

Und von dem her,

01:01:17 Sprecher 2

Die Schritte nachher waren eigentlich eher einfach Folgen.

01:01:21 Sprecher 2

Und der Schritt in die USA war schon ein großer Schritt.

01:01:23 Sprecher 2

Weil damals eben ohne Internet und so weiter wusste es und nicht
wirklich, wie das funktioniert und ob es dir gefällt.

01:01:32 Sprecher 2

Und ich glaube, das, was ich .

01:01:38 Sprecher 2

am meisten sagen würde, ist meine Voraussage, was in den nächsten .

01:01:44 Sprecher 2

wo ich und was ich in den nächsten zwei Jahren .

01:01:46 Sprecher 2

in zwei Jahren machen wäre, ist eigentlich oft falsch gewesen.

01:01:51 Sprecher 2

Zum Beispiel als ich in die USA gegangen bin, habe ich eigentlich
gedacht, geht für vier Jahre und komm dann in die Schweiz zurück.

01:01:59 Sprecher 2

Das war so.

01:02:00 Sprecher 2

Wenn du mich gefragt hättest.

01:02:03 Sprecher 2

Und wenn du mich zwei Jahre vor Google gefragt hättest, ob ich
irgendwann im Internet arbeite, das sagt wahrscheinlich.

01:02:13 Sprecher 1

Nicht.

01:02:14 Sprecher 1

Jetzt noch mal .

01:02:16 Sprecher 1

Einfach zum Abschluss, gibt es etwas vielleicht, was ich nicht gefragt
habe, was du aber findest, es wäre wichtig für deine Geschichte.

01:02:27 Sprecher 2

Jein, also ich glaube, etwas, was du hast ja gefragt über Management und
so weiter.

01:02:34 Sprecher 2

Ich glaube, was man unterschätzt.

01:02:37 Sprecher 2

Manchmal fragen mich die Leute, okay, wenn du jetzt 20 Jahre zurückgehen
könntest und sozusagen dir einen Rat geben könntest, nicht
professionally.

01:02:47 Sprecher 2

Was würdest du sagen?

01:02:48 Sprecher 2

Und ich würde sagen, dass, auch wenn man sehr technische Sachen macht,
dass der menschliche Aspekt unglaublich wichtig ist.

01:03:00 Sprecher 2

Das heißt, in den USA, sagen wir in Culture, also Work Culture oder usw.

01:03:08 Sprecher 2

Es gibt zum Beispiel den Spruch "Culture eats strategy for breakfast".

01:03:13 Sprecher 2

Also die Kultur ist die Strategie als Frühstück.

01:03:19 Sprecher 2

Und was es bedeuten will, ist, wenn du ein Team hast oder eine Firma,
die hat die falsche Kultur, dann spielt es keine Rolle, ob du die
richtige Strategie hast oder nicht.

01:03:28 Sprecher 2

Es wird nicht funktionieren.

01:03:31 Sprecher 2

Und wenn du die richtige Kultur hast, musst du natürlich auch noch eine
gute Strategie haben.

01:03:37 Sprecher 2

Und ich glaube, als Techniker, ich sah mich wirklich als reiner
Techniker, ist mir

01:03:44 Sprecher 2

wurde mir dann schon bewusst, wie wichtig das ist.

01:03:46 Sprecher 2

Aber vielleicht eben, würde ich sagen, als Story und eben als Ratschlag,
auch wenn man Informatiker ist, auch wenn man Technik ist, Techniker
ist, auch wenn man eine technische Gruppe leitet, ist es eben sehr
wichtig, eine Umgebung zu schaffen, wo sich die Leute wohlfühlen und
insbesondere, wo sie sich

01:04:08 Sprecher 2

über Probleme streiten können, sozusagen, ohne dass der Streit etwas
Negatives ist.

01:04:13 Sprecher 2

Dass man sozusagen nicht die Person angreift, sondern das Thema.

01:04:18 Sprecher 2

Und dass man dann eben auch Lösungen finden kann.

01:04:21 Sprecher 2

Oder zum Beispiel, wenn man eine Audit hat, das eben nachher als
Learning Opportunity brauchen kann und nicht als Blaming Opportunity.

01:04:33 Sprecher 2

Nicht, dessen Fehler war das.

01:04:34 Sprecher 2

Und das glaube ich, dass das

01:04:37 Sprecher 2

würde ich sagen, das hätte ich nicht gewusst beim Studium, aber nach dem
Studium nicht, wie wichtig das ist.

01:04:46 Sprecher 1

Okay.

01:04:49 Sprecher 1

Vielen Dank.

01:04:50 Sprecher 1

Ich glaube, unsere Zeit ist aufgebraucht.

01:04:55 Sprecher 2

Bitte, danke.

01:04:57 Sprecher 1

Wird vielleicht mal einen E-Mail-Austausch geben oder so in der
Zwischenzeit.

01:05:02 Sprecher 1

Aber es war jetzt sehr nett von Ihnen oder von dir.

01:05:07 Sprecher 1

Gut, danke vielmals.

01:05:10 Sprecher 2

Bitte, bitte.

01:05:12 Sprecher 2

So cool.

01:05:13 Sprecher 2

So now on time, yeah, on budget.
